\section{近期研究（2026 年 8--11 月）}
\paragraph{UAV3DCrop：分开评测作物航拍的外观、几何和冠层高度}
2026.08~\cite{2608_06404}。UAV3DCrop 把重复倾斜航拍的作物重建拆成逐场景优化和零样本前馈两条评测线，分别看新视图、摄影测量参考深度、米制尺度和作物冠层高度。评测保留真实尺度与对齐条件的差异，以免形状相似性掩盖米制误差。

论文发现方法排序随目标而变：Splatfacto-big 的外观表现领先，Scaffold-GS 在深度及冠层高度上更强。当前的新视图划分主要测试插值，深度参考本身依赖摄影测量，在重复纹理、缺少纹理或随风运动的作物上不等同于独立测量真值。

\paragraph{STARS-GS：按结构拆分航拍区域并约束跨区表面}
2026.09~\cite{2609_03447}。STARS-GS 从多视图影像和 SfM 点云出发，按场景结构划区并调整相邻区域边界。相比逐个原语作约束，它把邻近高斯的相对空间组织加入优化，并按照局部几何特征调节表面正则，使分块表面更连续。

GauU-Scene 与 AIR-LONGYAN 用于主要定量和定性比较，UrbanScene3D 与 Mill19 用于补充定性分析。方法研究的是既有航拍数据的表面重建，不直接处理在线跟踪。

\paragraph{SkyAnchor：让无位姿地面序列更新已有米制航拍场景}
2026.09~\cite{2609_13903}。SkyAnchor 将已有米制航空 SfM 和航空高斯作为支架。系统先验证来自地面短序列的候选定位，得到稀疏全局锚点，再在带锚点约束的地面子图中传播尺度和位姿。地面高斯加入时保留航空表示，最后只做有限联合细化，因此不必重新联合求解全部航空与地面相机。

原文在七个有航空和地面采集的场景比较，HorizonGS 场景的连接图像按协议移除。低跨视角重叠区域仍可能产生小的残余错位，通用稠密匹配器的表面覆盖也不能保证米制对齐。结论适用于已经拥有航空几何和高斯支架的更新任务。

\paragraph{HLC-GS：用高度层风险图约束卫星地表模型}
2026.09~\cite{2609_16772}。卫星高斯场在屋顶边界可能同时投影出屋顶和地面，若对高度直接作透明度加权，结果可能落在实际没有表面的中间层。HLC-GS 以高度分散度和主层可靠性生成风险图，只在高风险像素校正不稳的主层，并抑制不一致的远侧或次层响应。

论文在四个 JAX 与三个 IARPA 区域评测有效非森林栅格。相对 EOGS，平均高程 MAE 从 1.46 米降至 1.18 米，RMSE 从 2.78 米降至 2.58 米，2.5 米容差内的比例由 86.09\% 升至 88.61\%。风险图只能约束已有原语，严重遮挡和欠观测下仍难恢复，植被和跨时相变化还可能引入其他误差。

\paragraph{VDGS：按可见性划分航拍场景并补偿增密}
2026.09~\cite{2609_23049}。分块航拍高斯训练中，空间邻近的区域可能有完全不同的相机覆盖。VDGS 从粗模型提取平面原语，结合可见性划分场景，以每个块在图像上的重渲染覆盖决定相机归属。训练时累计低可见锚点的梯度并作补偿，使弱观测位置仍有机会增加原语。

实验覆盖 UrbanScene3D、Mill19 和 MatrixCity 的五个场景。视角不均匀的 Residence 与 Sci-Art 更能体现方法收益，覆盖较均匀的场景中表现主要是有竞争力。论文使用留出相机的图像指标，这与测绘精度是不同的评价目标。

\paragraph{OpenFlyScan：先估计局部重建质量，再补采薄弱区域}
2026.09~\cite{2609_24253}。OpenFlyScan 用图像特征、前馈几何和相机关系预测各区域的高斯重建质量，在正式训练前找薄弱区域并生成补采任务。初始影像和补采影像一起训练高斯，现场反馈由远程工作站计算。

论文在独立 UAVFF3D 实采序列及 Expo West、滴水湖现场数据检查预测和补采效果。该系统不要求无人机配新增硬件，但仍需要地面计算与数据传输。报告的图像质量增益应结合是否增加观测量来比较。

\paragraph{DIBR-GS：用逐步合成的深度伪视图补足遥感观测}
2026.09~\cite{2609_35612}。DIBR-GS 先把单目估计的深度与 SfM 稀疏点对齐，并融合多视图来初始化锚点。表示端加入来自参考图像的跨视角外观特征，训练端逐步把观测深度重投影到未观测视点，为高斯提供附加一致性约束。新的伪视图传递已有几何证据，并未增加真实采集。

作者在 LEVIR-NVS 上比较遥感新视图合成，观察到结构边界和细节的改进。框架仍需逐场景优化，伪视图质量还取决于单目深度和稀疏几何的对齐，不能直接转用为前馈重建。

\paragraph{SemLite-GS：按建筑实例分配航拍高斯训练和网格细化}
2026.11~\cite{2611_19773}。SemLite-GS 将正射影像中的建筑实例身份传播到倾斜视角，用相机和建筑实例的观测关系筛选训练任务。在局部区域并行优化目标高斯时，仍保留固定的周边高斯来解释遮挡。网格阶段将建筑平面约束写入符号距离场细化，尽量保持屋顶和立面结构。

论文使用六个无人机航拍场景，其中五个有 LiDAR 点云用于定量比较。对轻量网格模块单独比较时，固定了实例高斯输入，以隔离网格生成的影响。实验以已配准影像和相机参数为前提，量化结论不包含仅定性评价的第六个场景。
